\chapter{Investments}

\section{Three brokers, three formats}

The application imports statements from three brokers, and none of them export
the same thing.

\begin{table}[h]
\centering
\begin{tabular}{@{}p{3.2cm}p{2.4cm}p{7.9cm}@{}}
\toprule
\textbf{Broker} & \textbf{Format} & \textbf{What the file actually contains} \\
\midrule
MyInvestor & CSV & Cash movements of the account. Fund purchases appear as
negative amounts, redemptions as positive ones. \\
Kraken & CSV & Individual trades: pair, side, cost and volume. \\
Trade Republic & PDF & A statement in prose, one transaction per paragraph. \\
\bottomrule
\end{tabular}
\caption{What each export really is.}
\end{table}

\section{One shape for all of them}

Each parser returns the same two dataclasses, so nothing downstream needs to know
which broker produced the data.

\begin{lstlisting}[style=py]
@dataclass
class Position:
    asset: str
    isin: str | None = None
    quantity: Decimal | None = None
    invested: Decimal = Decimal("0")
    withdrawn: Decimal = Decimal("0")
    realised_pnl: Decimal = Decimal("0")
    looks_sold: bool = False
    meta: dict = field(default_factory=dict)

@dataclass
class ParseResult:
    broker: str
    positions: list[Position]
    warnings: list[str] = field(default_factory=list)
\end{lstlisting}

The \code{warnings} list is the error channel. A parser never raises on a bad
file; it returns an empty position list and an explanation. The upload page then
shows the explanation and the request completes normally. An exception would
produce a 500 page that tells the user nothing.

\begin{lstlisting}[style=py]
def parse(broker: str, content: bytes) -> ParseResult:
    module_path = PARSER_MODULES.get(broker)
    if module_path is None:
        return ParseResult(broker, [], [f"Unknown broker: {broker}"])
    return import_module(module_path).parse(content)
\end{lstlisting}

The dispatcher imports the parser module by name, at call time. That is not
laziness for its own sake: the parsers pull in pandas and pdfplumber, which take
seconds to import, and no user should pay that cost on a request that has
nothing to do with investments.

\section{Reading each format}

\subsection{MyInvestor}

The export is a cash statement, so the first job is telling fund movements apart
from transfers, interest and tax withholdings. Two rules do it: a list of words
that mark a cash movement, and the observation that fund names are printed in
upper case.

\begin{lstlisting}[style=py]
def is_fund_row(concept: str) -> bool:
    c = str(concept).strip()
    if any(w in c for w in SKIP_WORDS):
        return False
    if re.search(r"[a-z]", c):     # fund names are upper case
        return False
    return True
\end{lstlisting}

The remaining rows are grouped by fund. Negative amounts add up to
\code{invested}, positive ones to \code{withdrawn}. When the two nearly match,
the position is flagged as probably sold, and the preview shows the flag already
ticked.

\subsection{Kraken}

The export lists trades, so a position has to be reconstructed. For each pair,
the parser subtracts sells from buys, in both money and volume.

\begin{lstlisting}[style=py]
net_invested = buys["cost"].sum() - sells["cost"].sum()
volume = buys["vol"].sum() - sells["vol"].sum()
if volume > 1e-7 or abs(net_invested) > 0.01:
    positions.append(Position(asset=str(pair), ...))
\end{lstlisting}

The threshold on the volume avoids keeping positions that are only rounding
dust after everything was sold, while the money condition keeps a pair that is
closed but still shows a profit or a loss.

\subsection{Trade Republic}

The hardest one. The statement is a PDF written for humans, so the parser
flattens every page to text, splits it by transaction date, and matches each
block with regular expressions.

\begin{lstlisting}[style=py]
ISIN_RE = r"[A-Z]{2}[A-Z0-9]{9}[0-9]"

std = re.search(rf"(Buy|Sell|Savings).*?({ISIN_RE}).*?quantity:\s*([\d\.]+)",
                desc, re.IGNORECASE)
\end{lstlisting}

The extraction is fragile by nature; a change in the statement layout will break
it. Three details make it less fragile than it could be. Header and footer lines
are removed first. Amounts glued to quantities by the PDF text extraction are
split again with targeted substitutions. And the ISIN is matched by its actual
pattern, two country letters, nine alphanumerics and a check digit, rather than
by a hardcoded list of the products the author happens to own.

\begin{keypoint}[Generic beats specific]
The original private version had a special case for one particular product, with
its ISIN written in the source. The public version replaced it with a second,
looser regular expression that accepts any ISIN followed by a quantity. The
parser became shorter and covers more cases.
\end{keypoint}

Trade Republic also allows partial sells, so the parser keeps a running average
cost per ISIN and realises the profit of each sell:

\begin{lstlisting}[style=py]
elif "sell" in op:
    if qty > 1e-6:
        avg = cost / qty
        sold_cost = avg * q
        pnl += total - sold_cost
        qty -= q
        cost -= sold_cost
\end{lstlisting}

That realised profit travels in \code{Position.realised\_pnl} and is added back
later, because otherwise a profit that was already taken would disappear from
the totals.

\section{Profit depends on the broker}

Given the three formats above, a single profit formula would be wrong for at
least two of them. The calculation branches instead:

\begin{lstlisting}[style=py]
def total_profit(broker, invested, withdrawn, current_value,
                 realised_pnl=ZERO, sold=False) -> Decimal:
    if broker == "myinvestor":
        val_calc = ZERO if (sold and withdrawn > invested * HALF) else current_value
        return ((val_calc + withdrawn) - invested).quantize(Decimal("0.01"))
    if broker == "kraken":
        return (current_value - invested).quantize(Decimal("0.01"))
    if broker == "traderepublic":
        return ((current_value - invested) + realised_pnl).quantize(Decimal("0.01"))
    return ((current_value + withdrawn) - invested + realised_pnl).quantize(Decimal("0.01"))
\end{lstlisting}

Each branch answers a different question:

\begin{itemize}
  \item \textbf{MyInvestor} reports redemptions, so money taken out counts. When
        a fund was sold and most of the money already came back, the current
        value is ignored, otherwise the same money would be counted twice.
  \item \textbf{Kraken} nets buys against sells during parsing, so the profit is
        simply what the position is worth now minus what is still committed.
  \item \textbf{Trade Republic} already realised part of the profit on partial
        sells, so it is added to the unrealised part.
\end{itemize}

Every branch is covered by a test, including the double counting case, which is
the one that is easy to get wrong and impossible to notice by eye.

\section{Valuation: manual first}

A position needs a current value. Fetching it automatically works well for
listed products and badly for funds, whose share classes are hard to identify
and whose prices are often stale or missing.

The application therefore makes automatic pricing \textbf{opt-in}. When adding a
position the user ticks a box, or pastes a symbol; otherwise the position keeps
the value that was typed and only the profit is recomputed.

\begin{lstlisting}[style=py]
def refresh_position(inv: Investment) -> bool:
    if not inv.auto_value:
        return False
    symbol = inv.yahoo_symbol or (
        crypto_symbol(inv.asset) if inv.broker == "kraken" else None)
    if not symbol or inv.quantity is None:
        return False
    price = fetch_price(symbol)
    if price is None:
        return False
    inv.current_price = price
    inv.price_updated_at = datetime.now(timezone.utc)
    inv.current_value = (Decimal(inv.quantity) * price).quantize(Decimal("0.01"))
    inv.profit = total_profit(...)
    return True
\end{lstlisting}

The function returns a boolean, and every guard returns \code{False} instead of
raising. A failed price lookup is normal: markets close, tickers change, networks
fail. The position simply keeps its previous value, which is stale but correct,
rather than being wiped to zero.

\begin{pitfall}[Silent failure is right here, wrong elsewhere]
Swallowing an error is usually a bad habit. It is right here because the
alternative, an exception during a nightly job, would leave the whole portfolio
unrefreshed because of one delisted ticker. Compare with the backup in chapter
\ref{ch:operations}, where a failure must be loud.
\end{pitfall}

The batch version collects errors per position and keeps going:

\begin{lstlisting}[style=py]
for inv in positions:
    try:
        if refresh_position(inv):
            updated += 1
    except Exception as exc:      # one bad symbol must not abort the batch
        errors.append(f"{inv.asset}: {exc}")
\end{lstlisting}

\section{From an ISIN to a ticker}

Users know the ISIN of a product because it is printed on every statement. They
rarely know its Yahoo Finance symbol, which also depends on the exchange. The
application resolves one into the other in two steps.

\begin{enumerate}
  \item \textbf{Yahoo search.} Every suggested symbol is probed with a price
        request, preferring symbols quoted in euros. This is what resolves
        funds, whose symbols look like \code{0P0001CLDK.F}.
  \item \textbf{OpenFIGI.} A free mapping API that returns a ticker and an
        exchange code, which are combined into a Yahoo symbol through a suffix
        table. This works well for listed shares and ETFs.
\end{enumerate}

\begin{lstlisting}[style=py]
EXCH_TO_YAHOO_SUFFIX = {
    "GR": ".DE", "GY": ".DE", "GF": ".DE",   # Germany
    "SM": ".MC",                              # Spain
    "FP": ".PA", "NA": ".AS", "BB": ".BR",    # Euronext
    "LN": ".L", "IM": ".MI", "SW": ".SW",
    "US": "", "UN": "", "UW": "", "UQ": "",   # United States, no suffix
}
\end{lstlisting}

When a product is listed on several exchanges, European venues are preferred, so
the price comes back in the same currency as the rest of the portfolio.

Results are cached in a module-level dictionary, including the negative ones. An
ISIN that resolves to nothing resolves to nothing all day, and there is no reason
to ask twice.

\section{Contributions}

Regular investing means adding money to the same position every month. The model
keeps a single \code{invested} figure, the accumulated cost, which grows with
every contribution. Profit is then always the value minus what was actually put
in.

\begin{lstlisting}[style=py]
def contribute(db, inv, amount, units=ZERO) -> None:
    inv.invested = Decimal(inv.invested or 0) + amount
    if units:
        inv.quantity = Decimal(inv.quantity or 0) + Decimal(units)

    if not (inv.auto_value and refresh_position(inv)):
        inv.current_value = Decimal(inv.current_value or 0) + amount
        inv.profit = total_profit(...)
\end{lstlisting}

The conditional is the interesting part. For an automatically valued position,
the market price already reflects the new units, so adding the contributed amount
to the value would count it twice. For a manual position there is no price, so
the contributed money enters the value directly and the user adjusts later.

A position can also carry a \code{monthly\_contribution}. A scheduled job on the
first of each month applies all of them, logging and skipping any that fail, so
one broken position does not stop the rest.
